\documentclass[10pt]{article}
\usepackage{amsmath, amssymb, amsthm}
\usepackage{geometry}
\geometry{a4paper, margin=1in}

\usepackage{lmodern}
\usepackage{parskip}
\usepackage{enumitem}
\usepackage{hyperref}
\usepackage{multicol}
\usepackage{xcolor}
\usepackage{fancyhdr}
\usepackage{graphicx}

\definecolor{mypink}{RGB}{255, 105, 180}
\definecolor{mydarkpink}{RGB}{199, 21, 133}
\definecolor{mygray}{RGB}{80, 80, 80}
\hypersetup{
    colorlinks=true,
    linkcolor=mydarkpink,
    urlcolor=mydarkpink,
    citecolor=mydarkpink
}


\def\HS{\hspace{\fontdimen2\font}}

\pagestyle{fancy}
\fancyhf{}
\fancyhead[L]{\textcolor{mygray}{Elijah's Math Notes}}
\fancyhead[R]{\textcolor{mygray}{\thepage}}
\fancyfoot[C]{\textcolor{mygray}{Prepared by Elijah}}

\usepackage{titling}
\pretitle{\begin{center}\Large\color{mypink}}
\posttitle{\par\end{center}\vskip 0.5em}
\preauthor{\begin{center}\large\color{mygray}}
\postauthor{\par\end{center}}
\predate{\begin{center}\small\color{mygray}}
\postdate{\par\end{center}}

\theoremstyle{plain}
\newtheorem{theorem}{Theorem}[section]
\newtheorem{lemma}[theorem]{Lemma}
\newtheorem{proposition}[theorem]{Proposition}
\newtheorem{corollary}[theorem]{Corollary}

\theoremstyle{definition}
\newtheorem{definition}[theorem]{Definition}
\newtheorem{example}[theorem]{Example}

\theoremstyle{remark}
\newtheorem{remark}[theorem]{Remark}
\newtheorem{note}[theorem]{Note}

\newcommand{\R}{\mathbb{R}}
\newcommand{\N}{\mathbb{N}}
\newcommand{\Z}{\mathbb{Z}}
\newcommand{\todo}[1]{\textcolor{mydarkpink}{\textbf{TODO:} #1}}
\newcommand{\highlight}[1]{\textcolor{mypink}{#1}}

\renewcommand{\theenumi}{\Alph{enumi}}

\title{Vermont State Mathematics Coalition Talent Search Test \#1}
\author{Elijah Renner}
\date{\today}

\begin{document}

\maketitle
\tableofcontents

\section*{A Note}

I've never written any proofs or done competition math. I hope to improve, and I'd appreciate if you have resources to make my proofs more convincing and succinct. I just thought I'd give this a try.

-Elijah

\section{Problem \#1}

\paragraph{Problem} Find the greatest possible value of the ratio \(\frac{ABC-DEF}{GH-IJ}\) where \(A,B,C,D,E,F,G,H,I,J\) are distinct digits.

\paragraph{Solution} The problem requires us to maximize the ratio \(\frac{ABC-DEF}{GH-IJ}\). To do so, we must maximize the numerator, \(ABC-DEF\), and minimize the denominator, \(GH-IJ\).
Consider the case where \(GH-IJ=1\), i.e., the denominator cannot be further minimized. There are seven possible valid assignments such that \(GH-IJ=1\):\

\begin{enumerate}
\item \(GH=20, IJ=19\)
\item \(GH=30, IJ=29\)
\item \(GH=40, IJ=39\)
\item \(GH=50, IJ=49\)
\item \(GH=60, IJ=59\)
\item \(GH=70, IJ=69\)
\item \(GH=80, IJ=79\)	
\end{enumerate}

In all other assignments such that \(GH-IJ=1\), there are repeating digits.\footnote{To prove this assumption,  consider any assignment \(GH = XZ\) where \(Z \neq 0\). For \(GH-IJ=1\) to be true, \(XZ-(XZ-1)=1\) must also be true, it becomes clear that \(X\) always repeats. For the assignments \(GH=10, IJ=09\) and \(GH=90, IJ=89\), there are repeating digits.} Next, in order to ensure \(ABC-DEF\) is maximized, let \(E\) be the set of six remaining digits not assigned initially to \(GH\) or \(IJ\) in ascending order, and let \(E_n\) denote the \(n\)-th element in the set. Recall that to maximize the \(ABC-DEF\), we must maximize \(ABC\) and minimize \(DEF\). If our goal is to maximize \(ABC\), then \(A>B>C\). Similarly, to minimize \(DEF\), \(D<E<F\). So, using our list of the remaining digits \(E\), the assignment \(ABC - DEF = \overline{E_6E_5E_4} - \overline{E_1E_2E_3}\) maximizes \(ABC - DEF\). Using this procedure to assign digits to \(A,B,C,D,E,F\), we can simply test each of the seven possible assignments:\

\begin{enumerate}
\item \(GH=20, IJ=19 \implies \frac{ABC-DEF}{GH-IJ}=\frac{876-345}{1}=531\) 
\item \(GH=30, IJ=29 \implies \frac{ABC-DEF}{GH-IJ}=\frac{876-145}{1}=731\) 
\item \(GH=40, IJ=39 \implies \frac{ABC-DEF}{GH-IJ}=\frac{876-125}{1}=751\) 
\item \(GH=50, IJ=49 \implies \frac{ABC-DEF}{GH-IJ}=\frac{876-123}{1}=\fbox{753}\) 
\item \(GH=60, IJ=59 \implies \frac{ABC-DEF}{GH-IJ}=\frac{874-123}{1}=751\) 
\item \(GH=70, IJ=69 \implies \frac{ABC-DEF}{GH-IJ}=\frac{854-123}{1}=731\) 
\item \(GH=80, IJ=79 \implies \frac{ABC-DEF}{GH-IJ}=\frac{654-123}{1}=531\)	
\end{enumerate}

If \(753\) is the greatest possible value of the ratio when \(GH-IJ=1\), we must still consider the cases where \(GH-IJ \neq 1\). Consider the case where \(GH-IJ=2\). We want \(\frac{ABC-DEF}{GH-IJ}>753\). Inserting \(2\) as \(GH-IJ\) and manipulating, \(ABC-DEF>2(753)\). Since \(2(753)=1506\) and \(ABC-DEF\) can have at most three digits, it is not possible for \(\frac{ABC-DEF}{GH-IJ}>753\) when \(GH-IJ=2\). As \(GH-IJ\) becomes larger, \(ABC-DEF\) must also become larger, which it cannot given it can have at most three digits. Therefore, by contradiction, \(GH-IJ=1\) is optimal and \(\fbox{753}\) is the answer. 

\section{Problem \#2}

\subsection{a}

\paragraph{Problem} Ariel has 50\% more candy bars than Belle, who in turn has 50\% more candy bars than Cinderella. Ariel then gives 20 candy bars to Belle and \(N\) candy bars to Cinderella: afterwards, all three have the same number of candy bars. What is the value of \(N\)?

\paragraph{Solution} Let \(A,B,C\) represent the number of candy bars that Ariel, Belle, and Cinderella initially have, respectively. Since Ariel has 50\% more candy bars than Belle who has 50\% more candy bars than Cinderella, \(A=1.5\times 1.5C=2.25C\) and \(B=1.5C\). Applying the transformations,

\[A_{new}=2.25C-20-N\]
\[B_{new}=1.5C+20\]
\[C_{new}=C+N\]

It's given that \(A_{new}=B_{new}=C_{new}\), so 

\[2.25C-20-N=1.5C+20\]
\[1.5C+20=C+N\]

Manipulating the second equation to find \(C\) in terms of \(N\), \(C=2N-40\). Substituting \(2N-40\) for \(C\) into the first equation,

\[2.25(2N-40)-20-N=1.5(2N-40)+20\]

And, after solving for \(N\), the answer is \(\fbox{\(N=140\)}\).

\subsection{b}

\paragraph{Problem} Let \(A\) be the answer to part (a). A regular polygon has an internal angle measuring \(A+10\) degrees. How many diagonals does this polygon have?

\paragraph{Solution} Since \(A=140\), one interior angle of the regular polygon measures \(A+10=150\) degrees. The number of diagonals of an \(n\)-sided regular polygon is given by \(\frac{n(n-3)}{2}\). To find \(n\) given one interior angle measures \(150^{\circ}\), we can use the relationship between the interior angles and number of sides \(n\) of a regular polygon, \(\text{interior angle}_{\text{regular polygon}} = \frac{180(n - 2)}{n}\):

\[150^{\circ}=\frac{180(n-2)}{n}\implies n=12.\]

So, the polygon is a dodecagon. Now, by the the formula for the number of diagonals in a \(n\)-gon, \(\frac{n(n-3)}{2}\),

\[\text{\# diagonals = }\frac{n(n-3)}{2}=\frac{12(12-3)}{2}=54\]

Which means the answer is \(\fbox{54}\) diagonals.

\subsection{c}

\paragraph{Problem} Let \(B\) be the answer to part (b). A total of \(B\) blue marbles and 27 red marbles are placed into a bag. Two marbles are drawn randomly without replacement. What is the probability that the marbles are the same color?

\paragraph{Solution} Let \(R=27\) be the number of red marbles. Since \(B=54\), the total number of marbles \(B+R=81\). The probability of choosing 2 marbles of the same color is the sum of the probability of choosing two red marbles and the probability of choosing two blue marbles. The probability of choosing two red marbles is 

\[\frac{\text{ways to choose 2 red marbles from \(R\) red marbles}}{\text{ways to choose 2 marbles from \(R+B\) marbles}}=\frac{\binom{R}{2}}{\binom{R+B}{2}}=\frac{\binom{27}{2}}{\binom{81}{2}}=\frac{\frac{27!}{2!(27-2)!}}{\frac{81!}{2!(81-2)!}}=\frac{351}{3240}\]

Similarly, the probability of choosing two blue marbles is

\[\frac{\text{ways to choose 2 blue marbles from \(B\) blue marbles}}{\text{ways to choose 2 marbles from \(R+B\) marbles}}=\frac{\binom{B}{2}}{\binom{R+B}{2}}\frac{\binom{54}{2}}{\binom{81}{2}}=\frac{\frac{54!}{2!(54-2)!}}{\frac{81!}{2!(81-2)!}}=\frac{1431}{3240}\]

The sum of these two probabilities, i.e., the probability that two randomly chosen marbles are the same color, is \(\frac{351}{3240}+\frac{1431}{3240}=\boxed{\frac{11}{20}}\).

\section{Problem \#3}

\paragraph{Problem} Solve the cross-number puzzle below, where each entry is a digit from 1-9 (there are no 0s):

\begin{center}
\begin{tabular}{ |c|c|c|c| } 
 \hline
1 & 2 & 3 & 4 \\ 
\hline
5 &&& \\ 
\hline
6 &&& \\ 
 \hline
\end{tabular}
\end{center}

\begin{multicols}{2}

Across \\ 
1. A perfect cube. \\ 
5. A product of five distinct primes that sum to 42. \\ 
6. A multiple of 81 whose digits are in decreasing order. \\ 
Down \\
1. A prime less than 200. \\
2. An odd perfect square. \\
3. A multiple of 11. \\
4. A prime greater than 200.

\end{multicols}

\paragraph{Solution}

\begin{center}
\begin{tabular}{ |c|c|c|c| } 
 \hline
1 & 7 & 2 & 8 \\ 
\hline
9 & 2 & 8 & 2 \\ 
\hline
9 & 9 & 6 & 3 \\ 
 \hline
\end{tabular}
\end{center}\

Justification:

\textbf{1 Across}: The number is 1728.

- \(1728 = 12^3\), so it is a perfect cube.

\textbf{5 Across}: The number is 9282.

- Prime factorization: \(9282 = 2 \times 3 \times 7 \times 13 \times 17\).
- The primes are distinct and sum to \(2 + 3 + 7 + 13 + 17 = 42\).

\textbf{6 Across}: The number is 9963.

- The digits are in decreasing order: \(9 \geq 9 \geq 6 \geq 3\).
- \(9963 = 81 \times 123\), so it is a multiple of 81.


\textbf{1 Down}: The number is 199.

- 199 is a prime number less than 200.

\textbf{2 Down}: The number is 729.

- \(729 = 27^2\), which is an odd perfect square.

\textbf{3 Down}: The number is 286.

- \(286 = 11 \times 26\), so it is a multiple of 11.

\textbf{4 Down}: The number is 823.

- 823 is a prime number greater than 200.

This problem is best solved by listing the possible values for each row or column and eliminating them logically. First, I recognized that numbers across must be 4-digit numbers and numbers down must be 3-digit numbers. This allowed me to calculate the values (e.g., odd perfect squares) in the appropriate range. Then, you can eliminate possibilities with zeros. I also recognized that some digits must be odd, e.g., the last digit in the primes. At this point, I used the possible values' digits to eliminate numbers. For example, if I know that the possible second digits for 1 down are 2, 7, and 9, I can eliminate possibilities for 5 across whose first digit isn't 2, 7, or 9. Similar rules are applied until the problem is solved. 

\section{Problem \#4}

\paragraph{Problem} Find the unique ordered triple \(x,y,z\) of real numbers such that \(x^2+y^2+z^2+xy+xz+yz+1243=63x+47y+58z\).

\paragraph{Solution} First, collect terms on one side \[x^2+y^2+z^2+xy+xz+yz-63x-47y-58z+1243=0\] One approach is to eliminate the linear terms \(-63x-47y-58z\) to re-center the quadratic terms \(x^2+y^2+z^2+xy+xz+yz\) to the origin of some coordinate system defined by 
\[\begin{cases}
X=x-h\\
Y=y-k\\
Z=z-l
\end{cases}\iff \begin{cases}
x=X+h\\
y=Y+k\\
z=Z+l
\end{cases}\]

where \(h,k,l\) are constants. Substituting \(X,Y,Z\) for \(x,y,z\) into the original equation, 
\begin{align} 
(X+h)^2+(Y+k)^2+(Z+l)^2+(X+h)(Y+k)+(X+h)(Z+l) \\
+(Y+k)(Z+l)-63(X+h)-47(Y+k)-58(Z+l)+1243=0\nonumber 
\end{align}

Expanding reveals the quadratic terms \(X^2+Y^2+Z^2+XY+XZ+YZ\), linear terms \(X[2h+k+l-63]+Y[2k+h+l-47]+Z[2l+h+k-58]\), and constant terms \(h^2-63h+hk+hl+kl-47k+l^2-58l+1243+k^2\). To cancel the constant terms, 

\[\begin{cases}
2h+k+l=63\\
2k+h+l=47\\	
2l+h+k=58\\
\end{cases}\]

Solving this system, \((h,k,l)=(21,5,16)\). Substituting \((h,k,l)=(21,5,16)\) into the constant terms, \[(21)^2-63(21)+(21)(5)+(21)(16)+(5)(16)-47(5)+(16)^2-58(16)+1243+(5)^2=0\]

and, given that the linear terms cancel, since \(X^2+Y^2+Z^2+XY+XZ+YZ+X[2h+k+l-63]+Y[2k+h+l-47]+Z[2l+h+k-58]+h^2-63h+hk+hl+kl-47k+l^2-58l+1243+k^2=0\), \(X^2+Y^2+Z^2+XY+XZ+YZ\) must equal zero. Assuming the case where \(X=Y=Z=0\),  

\[\begin{cases}
0=x-h\\
0=y-k\\
0=z-l
\end{cases}\iff \begin{cases}
x=h\\
y=k\\
z=l
\end{cases}\]

So, we arrive at the ordered triple \((x,y,z)=(h,k,l)=(21,5,16)\). Inserting \((21,5,16)\) as \((x,y,z)\) into the original equation, we see that the triple \(\fbox{(21,5,16)}\) satisfies \[(21)^2+(5)^2+(16)^2+(21)(5)+(21)(16)+(5)(16)-63(21)-47(5)-58(16)+1243=0.\]

\section{Problem \#5}

\paragraph{Problem} Elias has a circular pizza. He randomly chooses 2025 pairs of points on the circumference of the pizza and makes a straight cut along the line joining each of his 2025 pairs of points. Find the expected number of pieces into which the pizza has been divided after Elias makes all 2025 cuts.

\paragraph{Solution} We will use Euler's formula for planar graphs, \(V-E+F=1\) where \(V\) is the number of vertices, \(E\) is the number of edges, and \(F\) is the number of faces, to answer this question by solving for \(F\).

First, we'll find \(V\). If 2025 pairs of points on a circle are randomly chosen and connected, \(N=2025\) chords will be randomly-placed on the circle. It follows that there will be \(\binom{N}{2}\) distinct pairs of chords. It is known that the probability of any two randomly-placed chords intersecting is \(\frac{1}{3}\), so the expected number of intersections is \(V=\text{probability of two randomly-placed chords intersecting} \times \text{pairs of chords}=\frac{1}{3}\binom{N}{2}\).

To find \(E\), we need to determine the expected number of edges per chord. If \(N\) chords are placed on a circle, it follows that any chord can intersect with \(N-1\) other chords. And, since we know the probability of any two random chords intersecting is \(\frac{1}{3}\), \(N_{\frac{\text{intersections}}{\text{chord}}}=\frac{N-1}{3}\). I'll now non-rigorously prove that the expected \(N_{\frac{\text{edges}}{\text{chord}}}=\frac{N-1}{3}+1\): if a chord is intersected 0 times, it will be divided into one edge. If a chord is intersected 1 time, it will be divided into two edges. If a chord is intersected twice, it will be divided into three edges, and so on. That is , each time an intersection is added to a chord, the number of edges will be equal to the number of intersections + 1.

So, \(E=N\times N_{\frac{\text{edges}}{\text{chord}}}=N(\frac{N-1}{3}+1)=\frac{N(N+2)}{3}.\)

We can now use the known values \(V=\frac{1}{3}\binom{N}{2}\text{ and }E=\frac{N(N+2)}{3}\) in the formula \(V-E+F=1\) to solve for \(F\), the number of faces or pizza pieces: \[\left(\frac{1}{3}\binom{N}{2}\right)-\left(\frac{N(N+2)}{3}\right)+F=1\implies \left(\frac{1}{3}\binom{2025}{2}\right)-\left(\frac{2025(2025+2)}{3}\right)+F=1\]

\[\implies F = \left(\frac{2025(2025+2)}{3}\right) - \left(\frac{1}{3}\binom{2025}{2}\right) + 1 = \fbox{685126} \text{ pizza pieces}\]

\section{Problem \#6}

\paragraph{Problem} Suppose that \(a\) and \(b\) are positive integers such that \(ab\) divides \(a^2+b^2+1\). Prove that both \(a\) and \(b\) must be Fibonacci numbers. (Recall that the Fibonacci numbers are defined by \(F_1=F_2=1\) and \(F_{n+1}=F_n+F_{n-1}\) for each \(n\geq2\).)

\paragraph{Proof} Since \(ab|a^2+b^2+1\), there exists some integer \(k\) such that \(a^2+b^2+1=kab\). Rewriting the equation as a quadratic in \(a\),\\ \begin{equation} \label{eq:1}a^2-kab+(b^2+1)=0\end{equation}\\ For integer solutions \(a\), the discriminant \(D=[kb]^2-4(b^2+1)=b^2(k^2-4)-4\) must be a perfect square. So, let \(D=s^2=b^2(k^2-4)-4\) for some integer \(s\geq0\). Our goal is to find integer solutions \((b,s)\) for \(k\geq3\) since for \(k=1,2\), the discriminant \(D\) becomes negative or zero, yielding non-integer solutions for \(a\), i.e., 
\[\text{For \(k=1\)}: s^2=(1^2-4)b^2-4=(-3)b^2-4<0\]
\[\text{For \(k=2\)}: s^2=(2^2-4)b^2-4=(0)b^2-4=-4\]

\textbf{Case 1: \(k=3\)}\

With \(k=3\), \[s^2=(3^2-4)b^2-4=5b^2-4\] which is a Pell's equation in the form \begin{equation} \label{eq:2}s^2-5b^2=-4 \end{equation} The fundamental solution is \((b_1,s_1)=(1,1)\): \[1^2-5(1)^2=-4\] It is known from the theory of Pell's equations that all such solutions (i.e., Pell multiples) \((b,s)\) are defined by a recurrence: \[b_{n+1}=9b_n+4s_n, \HS\HS\HS\HS\HS\HS\HS s_{n+1}=20b_n+9s_n\] \begin{align*}
(b_1, s_1) &= (1, 1), \\
(b_2, s_2) &= (9 \times 1 + 4 \times 1, 20 \times 1 + 9 \times 1) = (13, 29), \\
(b_3, s_3) &= (9 \times 13 + 4 \times 29, 20 \times 13 + 9 \times 29) = (233, 521), \\
&\vdots
\end{align*}
For the sake of brevity, we will assume this result and proceed with the proof. Notice that \(b\) corresponds to every other term in the Fibonacci sequence, i.e., the odd-indexed Fibonacci numbers. Next, it is known that some \(b\in \mathbb{Z}^+\) is a Fibonacci number \(\iff\) either \(5b^2+4\) or \(5b^2-4\) is a perfect square. Therefore, since \(s^2 \text{ (a perfect square)}=5b^2-4\), \(b\) must be a Fibonacci number. The solutions for equation \ref{eq:1} are \[a=\frac{kb\pm s}{2}=\frac{(3)b\pm s}{2}\]

And, since \(b\) and \(s\) are integers, and \(s\) is odd when \(b\) is odd (as in Fibonacci numbers at odd indices), 
\(3b\pm s\) is even, ensuring \(a\) is an integer. Computing \(a\) for the three \((b,s)\) generated by the recurrence, 

\[
\text{For \(b = 1, s = 1\):} \quad a = \frac{3 \times 1 \pm 1}{2} = \frac{3 \pm 1}{2} = \frac{4}{2}, \frac{2}{2} \implies a = 2, 1.
\]

\[
\text{For \(b = 13, s = 29\):} \quad a = \frac{3 \times 13 \pm 29}{2} = \frac{39 \pm 29}{2} = \frac{68}{2}, \frac{10}{2} \implies a = 34, 5.
\]

\[
\text{For \(b = 233, s = 521\):} \quad a = \frac{3 \times 233 \pm 521}{2} = \frac{699 \pm 521}{2} = \frac{1220}{2}, \frac{178}{2} \implies a = 610, 89.
\]

Observe that the values of \(a\) obtained are also Fibonacci numbers. This property holds for all valid solutions \((b,s)\) to the Pell's equation \ref{eq:2}, in which \(b\) is also an Fibonacci number.

\textbf{Case 2: \(k\geq4\)}\

We aim to show that for \(k \geq 4\), the equation
\[
s^2 = (k^2 - 4)b^2 - 4
\]
has no integer solutions \((b, s)\). Consider the equation in the form of the negative Pell's equation:
\[
s^2 - D b^2 = -4,
\]
where \(D = k^2 - 4\). For \(k \geq 4\), \(D \geq 12\). The negative Pell's equation \(x^2 - D y^2 = -4\) has integer solutions only for certain values of \(D\). In particular, for \(D > 5\), the equation \(x^2 - D y^2 = -4\) does not have integer solutions.

This can be seen by noting that the fundamental solution of the corresponding homogeneous Pell's equation \(x^2 - D y^2 = 1\) does not lead to solutions of the negative Pell's equation \(x^2 - D y^2 = -4\) when \(D > 5\). And, since \(D\) is not a perfect square and \(D > 5\), the minimal solution of \(x^2 - D y^2 = -4\) does not exist.

Therefore, for \(k \geq 4\), there are no integer solutions \((b, s)\) satisfying \(s^2 = (k^2 - 4)b^2 - 4\). Consequently, the discriminant \(D\) in equation \ref{eq:1} is not a perfect square, and \(a\) is not an integer.

Thus, the only feasible value of \(k\) such that \(a\) is an integer is \(k = 3\), leading to the equation \(s^2 = 5b^2 - 4\) for which integer solutions \((b, s)\) yield only Fibonacci numbers \(a\) and \(b\). Therefore, \(a\) and \(b\) must be Fibonacci numbers.


\end{document}

